% year: תשפ"ג
% type: מבחן
% semester: א
% moed: ב
% number: 3ב
% points: 10
% categories: func-analysis
% source: exams/מבחנים/תשפג/מועד ב תשפג.docx

\begin{question}
למשוואה $e^x=x+1$ הוכיחו כי קיים פתרון או הוכיחו כי לא קיים פתרון בתחום $x\neq0$.
\end{question}

\begin{hint}
חקרו את הפונקציה $g(x)=e^x-x-1$: מצאו את נקודת המינימום שלה.
\end{hint}

\begin{solution}
נוכיח ש\textbf{אין} פתרון עם $x\neq 0$. נגדיר $g(x)=e^x-x-1$, גזירה על $\R$, עם $g(0)=0$ ו-
\[ g'(x)=e^x-1. \]
עבור $x<0$: $g'(x)<0$, ולכן $g$ יורדת ממש ב-$(-\infty,0]$; עבור $x>0$: $g'(x)>0$, ולכן $g$ עולה ממש ב-$[0,\infty)$. (מונוטוניות ממש נובעת ממשפט לגרנז'.) לכן:
\begin{itemize}
  \item לכל $x<0$: $g(x)>g(0)=0$;
  \item לכל $x>0$: $g(x)>g(0)=0$.
\end{itemize}
כלומר $e^x>x+1$ לכל $x\neq0$, ולמשוואה $e^x=x+1$ אין פתרון בתחום $x\ne 0$ (הפתרון היחיד שלה הוא $x=0$). $\blacksquare$
\end{solution}
